% RESEARCH_DRAFT_SHA256: d534b82584cfe548cf60f0841a45b698c4d85ed0ee8c7d66053b2ba334d770f3
% Layout and source-reference wrappers expanded for the web renderer.
\documentclass[11pt]{article}
\usepackage[margin=25mm]{geometry}
\usepackage{amsthm}
\usepackage{fontspec,unicode-math}
\setmainfont{Latin Modern Roman}
\setmathfont{Latin Modern Math}
\usepackage{amsmath,amssymb,mathtools,xurl,hyperref}
\hypersetup{colorlinks=true,urlcolor=blue,linkcolor=blue}
\setcounter{secnumdepth}{0}
\setlength{\parindent}{0pt}
\setlength{\parskip}{5pt}
\setlength{\emergencystretch}{3em}
\newcommand{\sref}[2]{\hyperlink{source:#1:#2}{[S#2]}}
\newcommand{\cataloguescope}[1]{\paragraph{Source record.}#1}
\newcommand{\floor}[1]{\left\lfloor#1\right\rfloor}
\newtheorem{proposition}{Proposition}
\newtheorem{lemma}[proposition]{Lemma}
\begin{document}
\paragraph{Research manuscript.} OpenAI \texttt{gpt-6-astra}, reasoning \texttt{ultra}; 11 October 2026.
\paragraph{Preserved source body.} The complete catalogue mathematical text and sources follow with layout and source links expanded, including source qualifications and its historical preparation notes. New research begins at the section headed Research attempt.

% UnsolvedMath ID 30000196; source catalogue code OWR-792-010
\section{Observability for Polynomial Systems}\label{30000196}
{\small UnsolvedMath ID 30000196 · Computer Science}
\subsection{Definitions and mathematical statement}

\paragraph{Shared notation.}Standard number systems are $\mathbb N,\mathbb Z,\mathbb Q,\mathbb R,\mathbb C$. Specialized notation and index ranges follow the cited source and the definitions below.


A polynomial system has state and output dynamics given by polynomial maps. Observability concerns distinguishing initial states from outputs; algebraic observability uses algebraic information in the output and its derivatives.
\paragraph{Record scope and qualifications.}
The catalogue title is retained as metadata. Named certificates, contracts, bounds or frameworks in that title are not asserted as results of the original source extract.
The family of systems, coefficient topology, and precise meaning of generic must be supplied by the report.

Source-record limitation: The genericity family and algebraic-observability criterion are absent, while the definition-choice clause is a broad agenda.
\paragraph{Statement recorded in the supplied source.}
What is the right definition of observability for polynomial systems? Moreover, is the more restricted condition of algebraic observability generic? \href{https://doi.org/10.4171/owr/2005/11}{[S2]}
\paragraph{Source status.}
The supplied catalogue labels this source record partially solved; that metadata was checked against the supplied snapshot on 2026-10-11. The label is preserved for provenance and is not an independent verification of present open status or mathematical correctness. The catalogue literature assessment is dated 2026-08-21.
\subsection{Short English statement}
Identify an appropriate observability concept for polynomial systems and decide whether algebraic observability is generic.
\subsection{Sources}
\begin{itemize}
\item[S1] Ulam AI and the UnsolvedMath Contributors, UnsolvedMath: A Curated Collection of Open Mathematics Problems, record 30000196 (OWR-792-010). Catalogue curation and metadata: CC BY 4.0. The source extract is retained; definitions, English summaries and qualifications are editorial additions.
\url{https://huggingface.co/datasets/ulamai/UnsolvedMath}.
\item[S2] Regelungstheorie (2005), Oberwolfach Reports, DOI 10.4171/owr/2005/11, pp. 565--645
\url{https://doi.org/10.4171/owr/2005/11}.
\item[S3] generic algebraic-observability result
\url{https://www.researchgate.net/publication/228999815_Genericity_of_Algebraically_Observable_Polynomial_Systems}.
\item Source-specific Oberwolfach Reports licences, including CC BY-SA 4.0 where applicable, are retained; see THIRD\_PARTY\_NOTICES.md.
\end{itemize}
\paragraph{Preparation and review.}
The original source and UnsolvedMath contributors retain attribution for the source extract. Codex (GPT-6; exact serving revision not exposed) assisted with definitions, English summaries, scope notes and formatting. No new human mathematical authorship or independently verified solution is asserted. Permanent source ID 30000196 and code OWR-792-010 checked against the complete supplied registry and current repository. Distinct upstream records, including repeated questions and broad or damaged extracts, are retained for complete coverage of the source CS category. A repeated record is not a new independent conjecture. Complete mathematical-equivalence and current-literature review has not been performed. Preparation and final review follow contribution prompts C and D, with source gaps disclosed. Structural and rendering checks do not establish a new theorem or open status.
\label{end:30000196}

\section{Research attempt}
\subsection{Literature, source context and scope}
The full original catalogue body and source qualifications are preserved above. The relevant primary workshop passages were read in the cached publisher report. Historical claims and subsequent attributed results are distinguished from the mathematical work below. A restricted example, standard lemma or known theorem is not counted as new progress on the entire original target. Missing assumptions and repeated or rhetorical extracts are not silently turned into different solved problems.

\subsection{Distinguishability versus polynomial reconstruction}
We use the continuous-time Lie-derivative convention: for $\dot x=f(x)$ and $y=h(x)$, define $L_fh=h'f$ and iterate this operation. This is stated explicitly because the report also discusses discrete-time observations. The following one-dimensional examples test the inference that globally distinguishable real states must have polynomial reconstruction from finitely many output derivatives.
\begin{proposition}
The real polynomial system $\dot x=0,y=x^3$ distinguishes every pair of initial states but is not algebraically observable under this polynomial-reconstruction definition. Changing its dynamics to $\dot x=1$ while keeping $h=x^3$ makes it algebraically observable, with $x=y''/6$ along trajectories.
\end{proposition}
\begin{proof}
For the zero vector field, a trajectory starting at $x_0$ is constant, with output $x_0^3$. The cube map is injective on $\mathbb R$, so distinct initial states have different outputs already at time zero. All positive-order Lie derivatives vanish. Any proposed polynomial reconstruction from finitely many derivatives would therefore reduce to an identity $p(x^3)=x$ for a single-variable real polynomial $p$. A constant $p$ cannot work; if $p$ has positive degree $d$, the left side has degree $3d$, which cannot equal 1. Hence no such polynomial exists.
For $f=1$, $L_fh=3x^2$ and $L_f^2h=6x$, so the linear polynomial taking the third observation to one sixth its value reconstructs the state. Along a trajectory these Lie derivatives are the corresponding time derivatives of the output by the chain rule.
\end{proof}
This shows that the observable/algebraically observable distinction is substantive and that dynamics can supply missing algebraic information. The exceptional subfamily $f=0$ is not evidence against genericity in a larger coefficient space. The authored Helmke--Martin--Zhou preprint states generic algebraic observability for bounded positive degree bounds in a nonempty Zariski-open family, using continuous-time Lie derivatives; that is attributed literature, not a result of this attempt. Its full genericity proof is not reproved here.
\subsection{Verification and remaining gaps}
The Lie derivative chain for $f=1$ was checked exactly as $x^3,3x^2,6x,6,0$; for $f=0$ all derivatives after the first vanish. The degree obstruction is a complete proof for this example. No new genericity theorem, definition-selection criterion across all system classes, or effective general reconstruction algorithm was established. The source's mixed conventions are retained rather than silently equated.

\section{References and source attribution}
Primary workshop report: \url{https://doi.org/10.4171/owr/2005/11}. Relevant passages read: 2005/11 pp. 636--637, observability definitions and the genericity conjecture.

\url{https://www.researchgate.net/publication/228999815_Genericity_of_Algebraically_Observable_Polynomial_Systems}. Read the authored primary paper text hosted with the record, including the observation-algebra definition and Theorems 1.1--1.2. It states generic reconstruction for bounded positive degree families; the full proof was not independently verified here.

All retrievals were checked on 11 October 2026. Original question attribution remains in the preserved Sources section. External theorems are attributed and not claimed as new discoveries or fully verified merely by citation.


\paragraph{Integration review.} Codex (GPT-6; exact serving revision not exposed) reviewed the written arguments and expanded layout and source-reference wrappers. Repository metadata and the production web renderer were checked using Python 3.13 and JavaScript with MathJax 3.2.2. This model-assisted review is not independent human proof verification; no PDF compilation or formal proof checking is asserted.

\section{Completion Estimate}
\noindent\textbf{COMPLETION: 0\%}
No substantiated new progress toward the entire original target is claimed. The worked results and failed route are delimited above; a known result, illustrative model, or auxiliary obstruction is not counted as original-target completion.

\section{Final status and attribution}
\textbf{Attempt status: unresolved.} This describes this research attempt, not a new ruling on the source catalogue's status or on external literature claims.
The mathematics and this self-audit were generated by OpenAI \texttt{gpt-6-astra}, reasoning \texttt{ultra}, on 11 October 2026. Codex (GPT-6; exact serving revision not exposed) orchestrated the batch and contributed formatting/repository review. No human mathematical authorship, independent human proof checking, or formal verification is claimed. The historical preparation paragraphs in the reproduced catalogue body concern that body only.

Contribution prompt A was applied to the complete definition and supplied source record, followed by the adversarial checks of prompt B. Local reproduction files are \texttt{batch016/checks.py}, \texttt{computations.json}, and the per-ID audit JSON. They were not public links at preparation time. Numerical checks are finite and do not prove unrestricted results. No LaTeX engine was found locally; compilation was not performed by this research agent.

\end{document}
